% !TEX program = xelatex
% ============================================================
%  خطة تحويل «يَسِّر» إلى تطبيق PWA (Microservices-style)
%  الاستخدام على Overleaf:
%    1) أنشئ مشروعًا جديدًا فارغًا (Blank Project)
%    2) الصق محتوى هذا الملف في main.tex  (أو ارفع الملف)
%    3) Menu ← Compiler ← اختر: XeLaTeX   (مهم للعربية)
%    4) اضغط Recompile  ← ستحصل على PDF
% ============================================================
\documentclass[11pt]{article}

\usepackage{polyglossia}
\setmainlanguage{arabic}
\setotherlanguage{english}
\newfontfamily\arabicfont[Script=Arabic,Scale=1.0]{Amiri}
\newfontfamily\arabicfontsf[Script=Arabic,Scale=1.0]{Amiri}
\setmonofont{DejaVu Sans Mono}[Scale=0.85]

\usepackage{geometry}
\geometry{a4paper, margin=2.2cm}

\usepackage{xcolor}
\usepackage{graphicx}
\usepackage{booktabs}
\usepackage{longtable}
\usepackage{array}
\usepackage{enumitem}
\usepackage{titlesec}
\usepackage{fancyhdr}
\usepackage{hyperref}
\usepackage{tcolorbox}
\tcbuselibrary{skins,breakable}
\usepackage{fontawesome5}

% ── الألوان (هوية يَسِّر: ذهبي + داكن) ──
\definecolor{gold}{HTML}{B8933B}
\definecolor{darkbg}{HTML}{101414}
\definecolor{okgreen}{HTML}{27AE60}
\definecolor{warnred}{HTML}{C0392B}
\definecolor{softgray}{HTML}{F4F4F4}
\definecolor{midgray}{HTML}{7A7A7A}

% ── الترويسة والتذييل ──
\pagestyle{fancy}
\fancyhf{}
\renewcommand{\headrulewidth}{0.4pt}
\renewcommand{\footrulewidth}{0.4pt}
\fancyhead[R]{\small\color{gold}خطة تطبيق PWA — منصة يَسِّر}
\fancyhead[L]{\small\color{midgray}Microservices Roadmap}
\fancyfoot[C]{\small\thepage}

% ── عناوين الأقسام ──
\titleformat{\section}{\Large\bfseries\color{gold}}{\thesection}{0.6em}{}
\titleformat{\subsection}{\large\bfseries\color{darkbg}}{\thesubsection}{0.5em}{}
\titleformat{\subsubsection}{\normalsize\bfseries\color{midgray}}{\thesubsubsection}{0.4em}{}

\setlist[itemize]{rightmargin=0pt, itemsep=2pt, topsep=3pt}
\setlist[enumerate]{rightmargin=0pt, itemsep=2pt, topsep=3pt}

\hypersetup{
  colorlinks=true,
  linkcolor=gold,
  urlcolor=gold,
  pdftitle={خطة تطبيق PWA - منصة يَسِّر},
  pdfauthor={فريق يَسِّر}
}

% ── صندوق ملوّن ──
\newtcolorbox{infobox}[1][]{
  colback=softgray, colframe=gold, boxrule=0.6pt, arc=3pt,
  left=8pt, right=8pt, top=6pt, bottom=6pt, breakable, #1
}
\newtcolorbox{riskbox}[1][]{
  colback=red!4, colframe=warnred, boxrule=0.6pt, arc=3pt,
  left=8pt, right=8pt, top=6pt, bottom=6pt, breakable, #1
}
\newtcolorbox{okbox}[1][]{
  colback=green!4, colframe=okgreen, boxrule=0.6pt, arc=3pt,
  left=8pt, right=8pt, top=6pt, bottom=6pt, breakable, #1
}

% ── أوامر مساعدة ──
\newcommand{\code}[1]{\textcolor{darkbg}{\texttt{\small #1}}}
\newcommand{\ar}[1]{\textarabic{#1}}

\begin{document}

% ===================== الغلاف =====================
\begin{titlepage}
\centering
\vspace*{2.5cm}
{\Huge\bfseries\color{gold} يَسِّر}\\[0.4cm]
{\large\color{midgray} منصة تعليم الفقه الإسلامي}\\[2.5cm]

{\Huge\bfseries خطة تحويل الموقع إلى تطبيق}\\[0.6cm]
{\LARGE\bfseries PWA — بنية الخدمات المصغّرة}\\[1.2cm]

\rule{0.7\textwidth}{0.5pt}\\[0.8cm]
{\large Microservices-style Progressive Web App Roadmap}\\[0.8cm]
\rule{0.7\textwidth}{0.5pt}\\[2.5cm]

\begin{tabular}{r l}
\textbf{المنصة:}        & PWA (قابل للتثبيت — Android / iOS / Desktop) \\
\textbf{استراتيجية الكاش:} & شامل (App Shell + محتوى + أصول + خطوط) \\
\textbf{زر التحديث:}     & مُدمَج مع Service Worker \\
\textbf{المدة التقديرية:} & ٤.٥ أيام عمل \\
\textbf{النسخة:}         & 1.0 \\
\end{tabular}

\vfill
{\small\color{midgray} وثيقة تخطيط تقني — \today}
\end{titlepage}

\newpage
\tableofcontents
\newpage

% ===================== الملخص التنفيذي =====================
\section{الملخص التنفيذي}

\begin{infobox}
تحويل منصة «يَسِّر» من موقع ويب إلى \textbf{تطبيق قابل للتثبيت (PWA)} يعمل بلا إنترنت، مع الحفاظ الكامل على
\textbf{طبقة المزامنة السحابية الحالية} دون أي تغيير في منطقها. تُبنى الخطة على مبدأ \textbf{الخدمات المصغّرة}:
كل قدرة (تعريف، عمل دون اتصال، تحديث، تثبيت) في وحدة معزولة بواجهة واضحة، فلا يكسر تغييرٌ وحدةً أخرى.
\end{infobox}

\subsection*{الوضع الحالي}
\begin{tabular}{@{}p{5cm} p{9cm}@{}}
\toprule
\textbf{العنصر} & \textbf{الحالة} \\
\midrule
PWA & محذوف بالكامل (لا manifest ولا Service Worker فعّال) \\
المزامنة & \textcolor{okgreen}{تعمل فعلًا} — \code{onSnapshot} + حلّ تعارض بالطوابع الزمنية \\
العمل دون اتصال & كاش Firestore (IndexedDB 50MB) + كاش محتوى محلي \\
زر التحديث & فحص \code{index.html} يدويًا (بلا SW) \\
النشر & Firebase Hosting + تراويس CSP مُنفَّذة \\
\bottomrule
\end{tabular}

\subsection*{المبدأ المعماري}
\begin{center}
\renewcommand{\arraystretch}{1.35}
\begin{tabular}{@{}p{4cm} p{10cm}@{}}
\toprule
\textbf{الخدمة} & \textbf{المسؤولية} \\
\midrule
App Shell & هيكل React (موجود) \\
Content Service & تحميل المحتوى وكاشه (موجود) \\
\textbf{Sync Service} & \textbf{المزامنة السحابية (موجود — لا يُلمَس)} \\
Auth Service & المصادقة (موجود) \\
\addlinespace
Install Service & \textcolor{gold}جديد — واجهة التثبيت \\
Update Service & \textcolor{gold}جديد — دمج التحديث مع SW \\
Offline Service & \textcolor{gold}جديد — Workbox \\
Manifest Service & \textcolor{gold}جديد — التعريف والأيقونات \\
\bottomrule
\end{tabular}
\end{center}

% ===================== المراحل =====================
\section{مراحل التنفيذ}

\subsection{المرحلة ٠ — التحضير (نصف يوم)}
\begin{tabular}{@{}c p{6.5cm} p{5cm}@{}}
\toprule
\textbf{\#} & \textbf{المهمة} & \textbf{الملفات} \\
\midrule
٠.١ & إنشاء فرع \code{feat/pwa} & — \\
٠.٢ & إضافة \code{vite-plugin-pwa} كاعتماد تطوير & \code{package.json} \\
٠.٣ & توليد الأيقونات من الشعار الحالي & \code{build-brand-assets.mjs} \\
٠.٤ & نسخة احتياطية استباقية & \code{npm run backup} \\
٠.٥ & بوابة: الاختبارات خضراء قبل أي تغيير & — \\
\bottomrule
\end{tabular}
\begin{okbox}
\textbf{المخرج:} أيقونات جاهزة في \code{public/icons/}.
\end{okbox}

\subsection{المرحلة ١ — خدمة التعريف (Manifest) — نصف يوم}
الهدف: جعل الموقع قابلًا للتثبيت (يظهر خيار «تثبيت» في المتصفح).
\begin{tabular}{@{}c p{6.5cm} p{5cm}@{}}
\toprule
\textbf{\#} & \textbf{المهمة} & \textbf{التفصيل} \\
\midrule
١.١ & \code{public/manifest.webmanifest} & الاسم، \code{start\_url=/}, \code{display=standalone}, \code{dir=rtl}, \code{theme\_color}, \code{scope=/} \\
١.٢ & مصفوفة الأيقونات & 192، 512، maskable 512، apple-touch 180 \\
١.٣ & الربط في \code{index.html} & \code{link rel=manifest}, apple-touch-icon, theme-color \\
١.٤ & توسيع CSP & إضافة \code{manifest-src 'self'} \\
\bottomrule
\end{tabular}
\begin{okbox}
\textbf{بوابة القبول:} Lighthouse يعلن «Installable»، والـmanifest يُقرأ بلا أخطاء في DevTools.
\end{okbox}

\subsection{المرحلة ٢ — خدمة العمل دون اتصال (Workbox) — يوم}
هذه الخدمة معزولة تمامًا ولا تلمس التطبيق.
\begin{tabular}{@{}c p{6.5cm} p{5cm}@{}}
\toprule
\textbf{\#} & \textbf{المهمة} & \textbf{التفصيل} \\
\midrule
٢.١ & ضبط \code{vite.config.js} & \code{registerType=prompt}, \code{injectRegister=auto}, workbox \\
٢.٢ & تخزين App Shell مسبقًا & استثناء صريح لـ\code{lessons/**} و\code{books/**} \\
٢.٣ & كاش وقت التشغيل & stale-while-revalidate للأصول، NetworkFirst للخطوط، CacheFirst للصور \\
٢.٤ & استثناء Firestore & تجاهل \code{firestore.googleapis.com} (لا يمسّ IndexedDB) \\
٢.٥ & تعطيل التدمير الذاتي & إزالة \code{sw.js} و\code{service-worker.js} القديمين \\
٢.٦ & سلامة SEO & التأكد أن الصفحات الثابتة لا تتأثر بالـprecache \\
\bottomrule
\end{tabular}
\begin{riskbox}
\textbf{نقاط حرجة:}
\begin{itemize}
\item لا تمسّ كاش Firestore (يعمل في IndexedDB مستقل).
\item أخرج \code{lessons/**} من precache — وإلا انتفخ الـinstall.
\item خطوط Google: NetworkFirst (بلا نت ← آخر نسخة كاش).
\end{itemize}
\end{riskbox}
\begin{okbox}
\textbf{بوابة القبول:} اقطع النت ← التطبيق يفتح ويُقرأ المحتوى السابق؛ Console بلا أخطاء؛ المزامنة سليمة.
\end{okbox}

\subsection{المرحلة ٣ — خدمة التحديث (دمج SW) — نصف يوم}
مصدر حقيقة واحد للتحديث بدل فحص \code{index.html}.
\begin{tabular}{@{}c p{6.5cm} p{5cm}@{}}
\toprule
\textbf{\#} & \textbf{المهمة} & \textbf{التفصيل} \\
\midrule
٣.١ & استخدام \code{registerSW} & \code{onNeedRefresh}, \code{onOfflineReady} \\
٣.٢ & تعديل \code{useAppVersion} & \code{updateReady} من SW، \code{reloadNow} ← \code{updateSW(true)} \\
٣.٣ & إبقاء فحص \code{index.html} كـfallback & للمتصفحات القديمة \\
٣.٤ & البانر والإعدادات & ربط المصدر بلا تغيير الواجهة \\
٣.٥ & تحديث الاختبارات & \code{useAppVersion.test} + \code{UpdateBanner.test} \\
\bottomrule
\end{tabular}
\begin{riskbox}
أبقِ حارس \code{isFatalFirestoreError} و\code{reloadToLatest} في \code{firestoreRecovery.ts} — لا تلمسه.
\end{riskbox}
\begin{okbox}
\textbf{بوابة القبول:} انشر تعديلًا ← يظهر «نسخة جديدة» خلال دقائق، والزر يحدّث فعلًا.
\end{okbox}

\subsection{المرحلة ٤ — خدمة التثبيت (واجهة) — نصف يوم}
\begin{tabular}{@{}c p{6.5cm} p{5cm}@{}}
\toprule
\textbf{\#} & \textbf{المهمة} & \textbf{التفصيل} \\
\midrule
٤.١ & \code{usePWAInstall.ts} & التقاط \code{beforeinstallprompt}، كشف standalone، دعم iOS \\
٤.٢ & \code{InstallPrompt.tsx} & بانر أنيق بنمط التطبيق، زر تثبيت، تعليمات iOS \\
٤.٣ & الربط في \code{App.tsx} & يظهر مرة واحدة (خزّنه في localStorage) \\
٤.٤ & التكامل مع التعريف & دعوة التثبيت بعد شرائح التعريف \\
٤.٥ & احترام تقليل الحركة & كما في بقية المكونات \\
\bottomrule
\end{tabular}
\begin{okbox}
\textbf{بوابة القبول:} Android يعرض البانر ويثبّت؛ iOS يعرض النص التعليمي؛ لا يظهر إن كان مثبَّتًا.
\end{okbox}

\subsection{المرحلة ٥ — المزامنة: تحقق فقط (بلا تغيير كود) — ساعتان}
سياسة «ربط المزامنة عادي جدًا» — وهي فعليًا كذلك. هذه المرحلة تحقق لا بناء:
\begin{tabular}{@{}c p{10cm}@{}}
\toprule
\textbf{\#} & \textbf{الفحص} \\
\midrule
٥.١ & جهازان بنفس الحساب: تعديل في أحدهما يظهر في الآخر ≤ ٣٠٠ms \\
٥.٢ & وضع الطيران ثم تعديل ثم عودة: الدفع يُؤجَّل ثم يُرسل \\
٥.٣ & إغلاق التطبيق المثبَّت وفتحه: البيانات تستقر \\
٥.٤ & قفل الشاشة: الدفع الفوري يعمل \\
٥.٥ & مسح البيانات السحابية ثم «استئناف المزامنة» \\
٥.٦ & تأكيد أن SW لا يعترض \code{firestore.googleapis.com} \\
\bottomrule
\end{tabular}
\begin{infobox}
أي خلل يُعالَج في \code{useLibrarySync} بمعزلٍ تام، دون المساس بباقي الخدمات.
\end{infobox}

\subsection{المرحلة ٦ — الأمان والحوكمة — نصف يوم}
\begin{tabular}{@{}c p{6.5cm} p{5cm}@{}}
\toprule
\textbf{\#} & \textbf{المهمة} & \textbf{التفصيل} \\
\midrule
٦.١ & توسيع CSP & إضافة \code{manifest-src 'self'} و\code{worker-src 'self'} \\
٦.٢ & ترويسات SW & \textbf{إزالة} \code{Clear-Site-Data: "cache"} (يُفرّغ كاش PWA!) \\
٦.٣ & فحص SEO & التأكد من سلامة الصفحات الثابتة \\
٦.٤ & CI & فحص وجود manifest والأيقونات في البناء \\
٦.٥ & README & قسم «تطبيق مثبّت (PWA)» جديد \\
٦.٦ & تصحيح التوثيق & إزالة أي ذكر لـPWA «المحذوف» \\
\bottomrule
\end{tabular}
\begin{riskbox}
\textbf{مؤكد:} الترويسة الحالية على \code{/@(sw|service-worker).js} فيها \code{Clear-Site-Data: "cache"} — يجب إزالتها وإلا لا يعمل كاش الـPWA.
\end{riskbox}

\subsection{المرحلة ٧ — التحقق النهائي والنشر — يوم}
\begin{tabular}{@{}c p{5cm} p{7cm}@{}}
\toprule
\textbf{\#} & \textbf{البوابة} & \textbf{الأمر} \\
\midrule
٧.١ & فحص الأنواع & \code{npm run typecheck} \\
٧.٢ & ESLint & \code{npm run lint} \\
٧.٣ & الكود الميت & \code{npm run deadcode} \\
٧.٤ & الاختبارات & \code{npm test} \\
٧.٥ & البناء وSEO & \code{npm run build \&\& npm run check:seo} \\
٧.٦ & التكامل & \code{npm run test:e2e} \\
٧.٧ & Lighthouse & ≥ ٩٠، installable، offline يعمل \\
٧.٨ & نشر تجريبي & \code{npm run deploy:hosting} \\
٧.٩ & فحص حي على فون & تثبيت + تحديث + مزامنة + offline \\
\bottomrule
\end{tabular}

% ===================== المخاطر =====================
\section{المخاطر وسبل الوقاية}
\begin{center}
\renewcommand{\arraystretch}{1.3}
\begin{longtable}{@{}p{6cm} p{2cm} p{6.5cm}@{}}
\toprule
\textbf{الخطر} & \textbf{الاحتمال} & \textbf{الوقاية} \\
\midrule
\endfirsthead
\toprule
\textbf{الخطر} & \textbf{الاحتمال} & \textbf{الوقاية} \\
\midrule
\endhead
SW يخزّن Firestore فيكسر المزامنة & متوسط & قائمة منع + تجاهل النطاق صريحًا \\
\code{Clear-Site-Data: cache} يُلغي كاش PWA & \textcolor{warnred}{مؤكد} & إزالته من الترويسة (٦.٢) \\
تضخم precache (محتوى + PDF الكتاب) & متوسط & استثناء من patterns \\
CSP يحجب manifest/worker & مؤكد & توسيع \code{manifest-src}/\code{worker-src} \\
تنازع SW القديم (المُدمّر ذاتي) & منخفض & نسخ جديد يستبدله؛ منطق تنظيف انتقالي \\
زر التحديث يقول «محدّث» كاذبة & متوسط & مصدر SW واحد + اختبار \\
\bottomrule
\end{longtable}
\end{center}

% ===================== الجدول الزمني =====================
\section{الجدول الزمني التقديري}
\begin{center}
\renewcommand{\arraystretch}{1.3}
\begin{tabular}{@{}p{7cm} c@{}}
\toprule
\textbf{المرحلة} & \textbf{المدة} \\
\midrule
٠ — التحضير & ٠.٥ يوم \\
١ — التعريف (Manifest) & ٠.٥ يوم \\
٢ — العمل دون اتصال (Workbox) & ١ يوم \\
٣ — خدمة التحديث & ٠.٥ يوم \\
٤ — خدمة التثبيت & ٠.٥ يوم \\
٥ — تحقق المزامنة & ٠.٢٥ يوم \\
٦ — الأمان والحوكمة & ٠.٥ يوم \\
٧ — التحقق والنشر & ١ يوم \\
\midrule
\textbf{الإجمالي} & \textbf{≈ ٤.٥ أيام عمل} \\
\bottomrule
\end{tabular}
\end{center}

% ===================== معايير النجاح =====================
\section{معايير النجاح (Definition of Done)}
\begin{okbox}
\begin{itemize}
\item[\faCheck] الموقع قابل للتثبيت على Android / iOS / سطح المكتب.
\item[\faCheck] يفتح ويعمل بلا إنترنت بعد أول زيارة.
\item[\faCheck] زر «نسخة جديدة» يعمل عبر SW ويحدّث فعلًا.
\item[\faCheck] المزامنة بين الأجهزة تعمل كما هي (≤ ٣٠٠ms).
\item[\faCheck] كل البوابات خضراء (typecheck / lint / deadcode / test / build / seo / e2e).
\item[\faCheck] Lighthouse PWA ≥ ٩٠، وصفر أخطاء في Console.
\item[\faCheck] CSP محدّثة، والترويسات لا تُفرّغ كاش PWA.
\end{itemize}
\end{okbox}

% ===================== ملحق: هيكل الملفات =====================
\section{ملحق — هيكل الملفات بعد التنفيذ}
\begin{tcolorbox}[colback=darkbg, colframe=darkbg, coltext=white, arc=3pt, left=8pt, right=8pt]
\ttfamily\small\raggedleft
public/\\
\hspace*{1em}manifest.webmanifest\hspace{1em}\textcolor{gold}{جديد}\\
\hspace*{1em}icons/\hspace{1em}\textcolor{gold}{جديد}\\
\hspace*{2em}icon-192.png\\
\hspace*{2em}icon-512.png\\
\hspace*{2em}icon-maskable-512.png\\
\hspace*{2em}apple-touch-icon.png\\
src/\\
\hspace*{1em}hooks/\\
\hspace*{2em}usePWAInstall.ts\hspace{1em}\textcolor{gold}{جديد}\\
\hspace*{2em}useAppVersion.ts\hspace{1em}\textcolor{okgreen}{مُعدّل}\\
\hspace*{1em}Components/UI/\\
\hspace*{2em}InstallPrompt.tsx\hspace{1em}\textcolor{gold}{جديد}\\
vite.config.js\hspace{1em}\textcolor{okgreen}{مُعدّل — VitePWA}\\
firebase.json\hspace{1em}\textcolor{okgreen}{مُعدّل — CSP + headers}
\end{tcolorbox}

\vfill
\begin{center}
\rule{0.5\textwidth}{0.4pt}\\[0.4cm]
{\small\color{midgray} نهاية الوثيقة — منصة يَسِّر}
\end{center}

\end{document}
